\documentclass[../../main.tex]{subfiles}

\begin{document}

\section{Kakutani 不动点定理}\label{section:ch22-kakutani}

对两个群 $\mathcal{G}$ 和 $\mathcal{H}$, 映射 $\varphi:\mathcal{G}\to\mathcal{H}$ 称为群同态, 若对 $\mathcal{G}$ 中的每对元素 $g_{1}$, $g_{2}$, $\varphi(g_{1}\cdot g_{2})=\varphi(g_{1})\cdot\varphi(g_{2})$.

\begin{definition}[表示]\label{definition:ch22-representation}
令 $\mathcal{G}$ 为拓扑群而 $E$ 为 Banach 空间. 群同态 $\pi:\mathcal{G}\to\mathrm{GL}(E)$ 称为 $\mathcal{G}$ 在 $E$ 上的表示.
\end{definition}

与通常一样, 对 Banach 空间 $E$, 它的对偶空间, 即 $E$ 上的有界线性泛函空间, 记为 $E^{*}$. 回忆 $E^{*}$ 上的弱 $*$ 拓扑是使得对每个 $x\in X$, $E^{*}$ 上的泛函 $\psi\mapsto\psi(x)$ 连续的拓扑中具有最少集合个数的拓扑. Alaoglu 定理告诉我们 $E^{*}$ 的闭单位球关于弱 $*$ 拓扑是紧的.

\begin{definition}[伴随表示]\label{definition:ch22-adjoint-representation}
令 $\mathcal{G}$ 为拓扑群而 $E$ 为 Banach 空间, $\pi:\mathcal{G}\to\mathrm{GL}(E)$ 是 $\mathcal{G}$ 在 $E$ 上的表示. 伴随表示 $\pi^{*}:\mathcal{G}\to\mathrm{GL}(E^{*})$ 是 $\mathcal{G}$ 在 $E^{*}$ 上的表示, 它对 $g\in\mathcal{G}$ 定义为
\begin{align}
\text{对所有 }\psi\in E^{*},\quad\pi^{*}(g)\psi=\psi\circ\pi(g^{-1}).
\label{align:ch22-adjoint-def}
\end{align}
\end{definition}

我们把证明 $\pi^{*}$ 是一个群同态留作练习.

回忆向量空间 $V$ 上的度规或 Minkowski 泛函是一个正齐次的次可加泛函 $p:V\to\mathbf{R}$. 这样的泛函确定了局部凸拓扑向量空间 $V$ 的拓扑在原点的基. 在有了紧群 $\mathcal{G}$ 在 Banach 空间 $E$ 上的表示 $\pi$ 的条件下, 以下引理建立了 $E^{*}$ 上的一族由 $\mathcal{G}$ 参数化的正齐次次可加的泛函的存在性, 该族中的每个泛函在 $\pi^{*}$ 下不变且当限制在 $E^{*}$ 的有界子集上时, 关于弱 $*$ 拓扑是连续的.

\begin{lemma}[不变 Minkowski 泛函]\label{lemma:ch22-invariant-minkowski}
令 $\mathcal{G}$ 为紧群, $E$ 为 Banach 空间, 而 $\pi:\mathcal{G}\to\mathrm{GL}(E)$ 是 $\mathcal{G}$ 在 $E$ 上的表示. 令 $x_{0}$ 属于 $E$. 假设从 $\mathcal{G}$ 到 $E$ 的映射 $g\mapsto\pi(g)x_{0}$ 是连续的, 其中 $E$ 有范数拓扑. 定义 $p:E^{*}\to\mathbf{R}$ 为
\begin{align*}
\text{对 }\psi\in E^{*},\quad p(\psi)=\sup_{g\in\mathcal{G}}|\psi(\pi(g)x_{0})|.
\end{align*}
则 $p$ 是 $E^{*}$ 上的正齐次、次可加泛函. 它在 $\pi^{*}$ 下是不变的, 即
\begin{align*}
\text{对所有 }\psi\in E^{*}\text{ 与 }g\in\mathcal{G},\quad p(\pi^{*}(g)\psi)=p(\psi).
\end{align*}
此外, $p$ 在 $E^{*}$ 的任何有界子集上的限制关于 $E^{*}$ 上的弱 $*$ 拓扑是连续的.
\end{lemma}

\begin{proof}
由于 $\mathcal{G}$ 是紧的, 对 $\psi\in E^{*}$, 泛函 $g\mapsto\psi(\pi(g)x_{0})$ 在 $\mathcal{G}$ 上是连续的. $p:E^{*}\to\mathbf{R}$ 是恰当定义的. 显然 $p$ 是正齐次、次可加的, 且关于 $\pi^{*}$ 不变. 令 $B^{*}$ 为 $E^{*}$ 的有界子集. 为证明 $p:B^{*}\to\mathbf{R}$ 的弱 $*$ 连续性, 仅须证明对每个 $\psi_{0}\in B^{*}$ 和 $\varepsilon>0$, 存在 $\psi_{0}$ 的弱 $*$ 邻域 $\mathcal{N}(\psi_{0})$ 使得
\begin{align}
\text{对所有 }\psi\in\mathcal{N}(\psi_{0})\cap B^{*}\text{ 与 }g\in\mathcal{G},\quad|\psi(\pi(g)x_{0})-\psi_{0}(\pi(g)x_{0})|<\varepsilon.
\label{align:ch22-continuity-estimate}
\end{align}
令 $\psi_{0}$ 属于 $B^{*}$ 而 $\varepsilon>0$. 选取 $M>0$ 使得对所有 $\psi\in B^{*}$, $\|\psi\|\leqslant M$. 映射 $g\mapsto\pi(g)(x_{0})$ 是连续的而 $\mathcal{G}$ 是紧的. 因此存在 $\mathcal{G}$ 中的有限个点 $\{g_{1},\cdots,g_{n}\}$ 使得对每个 $k$, $1\leqslant k\leqslant n$, 存在 $g_{k}$ 的邻域 $\mathcal{O}_{k}$ 使得 $\{\mathcal{O}_{k}\}_{k=1}^{n}$ 覆盖 $\mathcal{G}$ 且对所有 $g\in\mathcal{O}_{g_{k}}$,
\begin{align}
\|\pi(g)x_{0}-\pi(g_{k})x_{0}\|<\varepsilon/4M.
\label{align:ch22-equicontinuity}
\end{align}
定义 $\psi_{0}$ 的弱 $*$ 邻域 $\mathcal{N}(\psi_{0})$ 为
\begin{align*}
\mathcal{N}(\psi_{0})=\{\psi\in E^{*}\mid(\psi-\psi_{0})(\pi(g_{k})x_{0})|<\varepsilon/2,\ 1\leqslant k\leqslant n\}.
\end{align*}
观察到对任何 $g\in\mathcal{G}$, $\psi\in E^{*}$, $1\leqslant k\leqslant n$,
\begin{align}
\psi(\pi(g)x_{0})-\psi_{0}(\pi(g)x_{0})=(\psi-\psi_{0})[\pi(g_{k})x_{0}]+(\psi-\psi_{0})[\pi(g)x_{0}-\pi(g_{k})x_{0}].
\label{align:ch22-minkowski-decomp}
\end{align}
为证明\eqref{align:ch22-continuity-estimate}, 令 $g$ 属于 $\mathcal{G}$ 而 $\psi$ 属于 $\mathcal{N}(\psi_{0})\cap B^{*}$. 选取 $k$, $1\leqslant k\leqslant n$, 使得 $g$ 属于 $\mathcal{O}_{k}$. 则 $|(\psi-\psi_{0})[\pi(g)x_{0}]|<\varepsilon/2$, 因为 $\psi$ 属于 $\mathcal{N}(\psi_{0})$. 另一方面, 由于 $\|\psi-\psi_{0}\|\leqslant 2M$, 我们从\eqref{align:ch22-equicontinuity}推出 $|(\psi-\psi_{0})[\pi(g)x_{0}-\pi(g_{k})x_{0}]|<\varepsilon/2$. 因此, 根据\eqref{align:ch22-minkowski-decomp}, \eqref{align:ch22-continuity-estimate}对 $\mathcal{N}(\psi_{0})$ 成立.
\end{proof}

\begin{definition}[不变集与不动点]\label{definition:ch22-invariant-set}
令 $\mathcal{G}$ 为拓扑群, $E$ 为 Banach 空间, 而 $\pi:\mathcal{G}\to\mathrm{GL}(E)$ 是 $\mathcal{G}$ 在 $E$ 上的表示. $E$ 的子集 $K$ 称为是在 $\pi$ 下不变的, 若对所有 $g\in\mathcal{G}$, $\pi(g)K\subseteq K$. 点 $x\in E$ 称为在 $\pi$ 下不动的, 若对所有 $g\in\mathcal{G}$, $\pi(g)x=x$.
\end{definition}

\begin{theorem}[Kakutani 型不动点]\label{theorem:ch22-kakutani-fixed-point}
令 $\mathcal{G}$ 为拓扑群, $E$ 为 Banach 空间, 而 $\pi:\mathcal{G}\to\mathrm{GL}(E)$ 是 $\mathcal{G}$ 在 $E$ 上的表示. 假设对每个 $x\in E$, 从 $\mathcal{G}$ 到 $E$ 的映射 $g\mapsto\pi(g)x$ 是连续的, 其中 $E$ 有范数拓扑. 假设存在 $E^{*}$ 的非空凸弱紧子集 $K^{*}$ 在 $\pi^{*}$ 下不变. 则存在 $K^{*}$ 中的泛函 $\psi$ 在 $\pi^{*}$ 下是不动的.
\end{theorem}

\begin{proof}
令 $\mathcal{F}$ 为 $K^{*}$ 的在 $\pi^{*}$ 下不变的非空、凸、弱 $*$ 闭子集合全体. 族 $\mathcal{F}$ 是非空的, 因为 $K^{*}$ 属于 $\mathcal{F}$. 用集合的包含关系对 $\mathcal{F}$ 排序. 这定义了 $\mathcal{F}$ 上的偏序. 由于 $K^{*}$ 属于 $\mathcal{F}$, $\mathcal{F}$ 是非空的. 每个全序子集有有限交性质. 但对任何紧拓扑空间, 具有有限交性质的非空闭子集族有非空交. 任何凸集族的交是凸的且任何 $\pi^{*}$ 不变集族的交是 $\pi^{*}$ 不变的. 因此 $\mathcal{F}$ 的每个全序子集以它的非空交作为下界. 我们从 Zorn 引理推出, 存在 $\mathcal{F}$ 中的集合 $K_{0}^{*}$ 关于包含关系是最小的, 即 $K_{0}^{*}$ 没有属于 $\mathcal{F}$ 的真子集. 这个最小子集是弱 $*$ 闭的, 因此是弱 $*$ 紧的. 我们重新标记且假设 $K^{*}$ 自身是这个最小子集.

我们宣称 $K^{*}$ 由单个泛函组成. 否则, 选取 $K^{*}$ 中的两个不同泛函 $\psi_{1}$ 与 $\psi_{2}$. 选取 $x_{0}\in E$ 使得 $\psi_{1}(x_{0})\neq\psi_{2}(x_{0})$. 定义泛函 $p:K^{*}\to\mathbf{R}$ 为
\begin{align*}
\text{对 }\psi\in K^{*},\quad p(\psi)=\sup_{g\in\mathcal{G}}|\psi(\pi(g)x_{0})|.
\end{align*}
由于 $K^{*}$ 是弱 $*$ 紧的, 一致有界原理告诉我们 $K^{*}$ 是有界的. 根据前一个引理, $p$ 关于弱 $*$ 拓扑是连续的. 因此, 若对 $r>0$ 与 $\eta\in K^{*}$, 我们定义
\begin{align}
B_{0}(\eta,r)=\{\psi\in K^{*}\mid p(\psi-\eta)<r\}\text{ 与 }\overline{B_{0}}(\eta,r)=\{\psi\in K^{*}\mid p(\psi-\eta)\leqslant r\},
\label{align:ch22-balls}
\end{align}
则 $B_{0}(\eta,r)$ 关于 $K^{*}$ 上的弱 $*$ 拓扑是开的, 而 $\overline{B_{0}}(\eta,r)$ 关于同一个拓扑是闭的. 这些集合的每个是凸的, 因为根据前一个引理, $p$ 是正齐次与次可加的.

定义 $d=\sup\{p(\psi-\varphi)\mid\psi,\varphi\in K^{*}\}$. 则 $d$ 是有限的, 由于 $p$ 在弱 $*$ 紧集 $K^{*}$ 上是连续的, 且 $d>0$, 因为 $p(\psi_{1}-\psi_{2})>0$. 由于 $K^{*}$ 是弱 $*$ 紧的且每个 $B_{0}(\eta,r)$ 是弱 $*$ 开的, 我们可以选取 $K^{*}$ 的有限子集 $\{\psi_{k}\}_{k=1}^{n}$ 使得
\begin{align*}
K^{*}=\bigcup_{k=1}^{n}B_{0}(\psi_{k},d/2).
\end{align*}
定义
\begin{align*}
\psi^{*}=\frac{\psi_{1}+\cdots+\psi_{k}+\cdots+\psi_{n}}{n}.
\end{align*}
由于 $K^{*}$ 是凸的, 泛函 $\psi^{*}$ 属于 $K^{*}$. 令 $\psi$ 为 $K^{*}$ 中的任何泛函. 根据 $d$ 的定义, 对 $1\leqslant k\leqslant n$, $p(\psi-\psi_{k})\leqslant d$. 由于 $\{B_{0}(\psi_{k},d/2)\}_{k=1}^{n}$ 覆盖 $K^{*}$, $\psi$ 属于某个 $B_{0}(\psi_{k_{0}},d/2)$. 因此, 根据 $p$ 的正齐次性与次可加性,
\begin{align*}
p(\psi-\psi^{*})\leqslant d',\text{ 其中 }d'=\frac{n-1}{n}d+\frac{d}{2}<d.
\end{align*}
定义
\begin{align*}
K'=\bigcap_{\psi\in K^{*}}\overline{B_{0}}(\psi,d').
\end{align*}
则 $K'$ 是 $K^{*}$ 的弱 $*$ 闭, 因此是 $K^{*}$ 的弱 $*$ 紧凸子集. 由于它包含泛函 $\psi^{*}$, 它是非空的. 我们宣称 $K'$ 在 $\pi^{*}$ 下是不动的. 为证实这一点, 对 $\eta\in K'$, $\psi\in K^{*}$ 与 $g\in\mathcal{G}$, 我们必须证明 $p(\pi^{*}(g)\eta-\psi)\leqslant d'$. 由于 $p$ 是 $\pi^{*}$ 不动的且 $p(\eta-\pi^{*}(g^{-1})\psi)\leqslant d'$,
\begin{align*}
p(\pi^{*}(g)\eta-\psi)=p(\eta-\pi^{*}(g^{-1})\psi)\leqslant d'.
\end{align*}
根据 $K^{*}$ 的最小性, $K^{*}=K'$. 这是一个矛盾, 由于根据 $d$ 的定义, 存在 $K^{*}$ 中的泛函 $\psi'$ 与 $\psi''$ 使得 $p(\psi'-\psi'')>d'$, 因此 $\psi''$ 不属于 $\overline{B_{0}}(\psi',d')$. 我们从这个矛盾推出 $K^{*}$ 由单个泛函组成. 证明完毕.
\end{proof}

\begin{definition}[正则表示]\label{definition:ch22-regular-rep}
令 $\mathcal{G}$ 为紧群, 而 $C(\mathcal{G})$ 为 $\mathcal{G}$ 上赋予最大值范数的连续实值函数的 Banach 空间. 谈到 $C(\mathcal{G})$ 上的正则表示 $\pi:\mathcal{G}\to\mathrm{GL}(C(\mathcal{G}))$, 我们意味着:
\begin{align*}
\text{对所有 }f\in C(\mathcal{G}),\ x\in\mathcal{G}\text{ 与 }g\in\mathcal{G},\quad[\pi(g)f](x)=f(g^{-1}\cdot x).
\end{align*}
\end{definition}

我们把证明正则表示的确是一个表示留作练习. 以下引理表明紧群 $\mathcal{G}$ 在 $C(\mathcal{G})$ 上的正则表示具有施加在定理\cref{theorem:ch22-kakutani-fixed-point}上的连续性.

\begin{lemma}[正则表示的连续性]\label{lemma:ch22-regular-continuity}
令 $\mathcal{G}$ 为紧群而 $\pi:\mathcal{G}\to\mathrm{GL}(C(\mathcal{G}))$ 是 $\mathcal{G}$ 到 $C(\mathcal{G})$ 上的正则表示. 则对每个 $f\in C(\mathcal{G})$, 从 $\mathcal{G}$ 到 $C(\mathcal{G})$ 的映射 $g\mapsto\pi(g)f$ 是连续的, 其中 $C(\mathcal{G})$ 有由最大值范数诱导的拓扑.
\end{lemma}

\begin{proof}
令 $f$ 属于 $C(\mathcal{G})$. 仅须检验映射 $g\mapsto\pi(g)f$ 在单位元 $e\in\mathcal{G}$ 是连续的. 令 $\varepsilon>0$. 我们宣称存在单位元的邻域 $\mathcal{U}$ 使得
\begin{align}
\text{对所有 }g\in\mathcal{U},\ x\in\mathcal{G},\quad|f(g\cdot x)-f(x)|<\varepsilon.
\label{align:ch22-uniform-cont}
\end{align}
令 $x$ 属于 $\mathcal{G}$. 选取 $x$ 的邻域 $\mathcal{O}_{x}$ 使得对所有 $x'\in\mathcal{O}_{x}$,
\begin{align*}
|f(x')-f(x)|<\varepsilon/2.
\end{align*}
因此,
\begin{align}
\text{对所有 }x',\ x''\in\mathcal{O}_{x},\ |f(x')-f(x'')|<\varepsilon.
\label{align:ch22-small-neigh}
\end{align}
根据群运算的连续性, 我们可以选取单位元的邻域 $\mathcal{U}_{x}$ 和 $x$ 的邻域 $\mathcal{V}_{x}$ 使得 $\mathcal{V}_{x}\subseteq\mathcal{O}_{x}$ 且 $\mathcal{U}_{x}\cdot\mathcal{V}_{x}\subseteq\mathcal{O}_{x}$. 根据 $\mathcal{G}$ 的紧性, 存在覆盖 $\mathcal{G}$ 的有限族 $\{\mathcal{V}_{x_{k}}\}_{k=1}^{n}$. 定义 $\mathcal{U}=\bigcap_{k=1}^{n}\mathcal{U}_{x_{k}}$. 则 $\mathcal{U}$ 是 $\mathcal{G}$ 中的单位元的邻域. 我们宣称\eqref{align:ch22-uniform-cont}对 $\mathcal{U}$ 的这个选取成立. 事实上, 令 $g$ 属于 $\mathcal{U}$ 而 $x$ 属于 $\mathcal{G}$. 则 $x$ 属于某个 $\mathcal{V}_{x_{k}}$. 因此,
\begin{align*}
x\in\mathcal{V}_{x_{k}}\subseteq\mathcal{O}_{x_{k}}\text{ 且 }g\cdot x\in\mathcal{U}\times\mathcal{V}_{x_{k}}\subseteq\mathcal{U}_{x_{k}}\times\mathcal{V}_{x_{k}}\subseteq\mathcal{O}_{x_{k}}.
\end{align*}
因此 $x$ 和 $g\cdot x$ 都属于 $\mathcal{O}_{x_{k}}$. 于是根据\eqref{align:ch22-small-neigh}, $|f(g\cdot x)-f(x)|<\varepsilon$. 因此\eqref{align:ch22-uniform-cont}得证. 用 $\mathcal{U}\cap\mathcal{U}^{-1}$ 代替 $\mathcal{U}$. 因此,
\begin{align*}
\text{对所有 }g\in\mathcal{U},\ x\in\mathcal{G},\quad|f(g^{-1}\cdot x)-f(x)|<\varepsilon.
\end{align*}
即
\begin{align*}
\text{对所有 }g\in\mathcal{U},\quad\|\pi(g)f-\pi(e)f\|_{\max}<\varepsilon.
\end{align*}
这证明了所要的连续性.
\end{proof}

对紧群 $\mathcal{G}$, 我们称泛函 $\psi\in[C(\mathcal{G})]^{*}$ 为概率泛函, 若它在常函数 $f=1$ 取值 $1$ 且在以下意义下是正的: 对 $f\in C(\mathcal{G})$, 若在 $\mathcal{G}$ 上 $f\geqslant 0$, 则 $\psi(f)\geqslant 0$.

\begin{theorem}[Kakutani 不动点定理]\label{theorem:ch22-kakutani}
令 $\mathcal{G}$ 为紧群而 $\pi:\mathcal{G}\to\mathrm{GL}(C(\mathcal{G}))$ 是 $\mathcal{G}$ 在 $C(\mathcal{G})$ 上的正则表示. 则存在概率泛函 $\psi\in[C(\mathcal{G})]^{*}$ 在伴随作用 $\pi^{*}$ 下不动, 即
\begin{align}
\text{对所有 }f\in C(\mathcal{G})\text{ 与 }g\in\mathcal{G},\quad\psi(f)=\psi(\pi(g)f).
\label{align:ch22-kakutani-invariant}
\end{align}
\end{theorem}

\begin{proof}
根据 Alaoglu 定理, $[C(\mathcal{G})]^{*}$ 的闭单位球是弱 $*$ 紧的. 令 $K^{*}$ 为 $C(\mathcal{G})$ 上的正概率泛函族. 观察到若 $\psi$ 是概率泛函而 $f$ 属于 $C(\mathcal{G})$ 满足 $\|f\|_{\max}\leqslant 1$, 则根据 $\psi$ 的正性和线性, 由于 $-1\leqslant f\leqslant 1$,
\begin{align*}
-1=\psi(-1)\leqslant\psi(f)\leqslant\psi(1)=1.
\end{align*}
因此 $|\psi(f)|\leqslant 1$, 且因此 $\|\psi\|\leqslant 1$. 因此 $K^{*}$ 是 $E^{*}$ 的闭单位球的凸子集. 我们宣称 $K^{*}$ 是弱 $*$ 闭的. 事实上, 对每个非负函数 $f\in C(\mathcal{G})$, 集合 $\{\psi\in[C(\mathcal{G})]^{*}\mid\psi(f)\geqslant 0\}$ 是弱 $*$ 闭的, 在常函数 $f\equiv 1$ 取值 $1$ 的泛函 $\psi$ 的集合也如此. 集合 $K^{*}$ 是弱 $*$ 闭集的交, 因此是弱 $*$ 闭的. 作为紧集的闭子集, $K^{*}$ 是弱 $*$ 紧的. 最后, 集合 $K^{*}$ 是非空的, 这是因为, 若 $x_{0}$ 是 $\mathcal{G}$ 中的任何点, 在属于 $K^{*}$ 的每个 $f\in C(\mathcal{G})$ 处取值 $f(x_{0})$ 的 Dirac 泛函属于 $K^{*}$.

显然 $K^{*}$ 在 $\pi^{*}$ 下是不变的. 前一个引理告诉我们正则表示具有应用定理\cref{theorem:ch22-kakutani-fixed-point}所需要的连续性. 根据该定理, 存在泛函 $\psi\in K^{*}$ 在 $\pi^{*}$ 下是不动的, 即\eqref{align:ch22-kakutani-invariant}成立.
\end{proof}

\begin{exercise}\label{exercise:ch22-13}
证明表示的伴随也是表示.
\end{exercise}

\begin{exercise}\label{exercise:ch22-14}
证明概率泛函的范数是 $1$.
\end{exercise}

\begin{exercise}\label{exercise:ch22-15}
令 $E$ 为自反 Banach 空间而 $K^{*}$ 是 $E^{*}$ 的凸子集, 关于由范数诱导的度量是闭的. 证明 $K^{*}$ 是弱 $*$ 闭的. 另一方面, 证明: 若 $E$ 不是自反的, 则 $E$ 的闭单位球在 $(E^{*})^{*}=E^{**}$ 的自然嵌入下的象是 $E^{**}$ 的凸闭子集, 它关于由范数诱导的度量是有界的但不是弱 $*$ 闭的.
\end{exercise}

\begin{exercise}\label{exercise:ch22-16}
令 $\mathcal{G}$ 为紧群, $E$ 为自反 Banach 空间, $\pi:\mathcal{G}\to\mathrm{GL}(E)$ 是表示. 假定对每个 $x\in E$, 映射 $g\mapsto\pi(g)x$ 是连续的. 假设存在 $E$ 的在 $\pi$ 下不变的非空强闭、有界、凸子集 $K$. 证明 $K$ 包含在 $\pi$ 下不动的点.
\end{exercise}

\begin{exercise}\label{exercise:ch22-17}
令 $\mathcal{G}$ 为拓扑群而 $E$ 为 Banach 空间, $\pi:\mathcal{G}\to\mathrm{GL}(E)$ 是表示. 对 $x\in E$, 证明映射 $g\mapsto\pi(g)x$ 是连续的当且仅当它在 $e$ 处是连续的.
\end{exercise}

\begin{exercise}\label{exercise:ch22-18}
假定 $\mathcal{G}$ 为拓扑群而 $X$ 为拓扑空间, 而 $\varphi:\mathcal{G}\times X\to X$ 是映射. 对 $g\in\mathcal{G}$, 定义映射 $\pi(g):X\to X$ 为对所有 $x\in X$, $\pi(g)(x)=\varphi(g,x)$. 为使得 $\pi$ 是 $\mathcal{G}$ 在 $C(X)$ 上的表示, $\varphi$ 必须具有什么性质? 为使得对每个 $x\in X$, 映射 $g\mapsto\pi(g)x$ 是连续的, $\varphi$ 必须具有什么进一步的性质?
\end{exercise}

\end{document}